\section{Hasil}

\subsection{E1: Overlay Dua Model pada Sumbu Bersama}

Gambar~\ref{fig:e1} menampilkan fraksi fermentatif versus ketersediaan
glukosa untuk ketiga organisme. Kedua model menghasilkan kurva yang nyaris
identik: pada glukosa rendah keduanya memprediksi respirasi murni
(fraksi fermentatif $\approx 0$), dan pada glukosa tinggi keduanya
memprediksi glikolisis dominan. Zona ketidaksetujuan (perbedaan
$> 0{,}1$) sangat sempit --- terbatas pada wilayah transisi.

\begin{figure}[t]
    \centering
    \includegraphics[width=\columnwidth]{F1_overlay.png}
    \caption{E1: Overlay kedua model pada sumbu glukosa bersama. Kedua model memprediksi fenotipe yang hampir identik. Area merah muda menandai zona ketidaksetujuan ($|\Delta| > 0{,}1$).}
    \label{fig:e1}
\end{figure}

\textbf{Temuan E1:} Kedua model \textit{tidak bertentangan secara struktural}
pada sumbu observabel bersama. Perbedaan bukan pada ``model mana yang benar'',
melainkan pada besaran efisiensi yang mereka hitung dari model yang sama.

\subsection{E2: Kapasitas versus Realisasi --- Titik-Balik Verdict}

Gambar~\ref{fig:e2} menunjukkan efisiensi terealisasi kedua jalur sebagai
fungsi fraksi pemanfaatan glikolitik $u_G$. Pada $u_G = 1$ (kapasitas
penuh, asumsi Kukurugya), glikolisis menang. Seiring $u_G$ turun
(enzim glikolitik idle, temuan Shen), respirasi menang.

Titik-balik terjadi tepat di $u_G = \rho$ (rasio kapasitas):

\begin{table}[h]
\centering
\caption{Titik-balik verdict $u_G$ dengan interval kepercayaan 90\%.}
\label{tab:flip}
\begin{tabular}{lccc}
\toprule
\textbf{Organisme} & $\boldsymbol{\rho}$ & \textbf{Flip} $\boldsymbol{u_G}$ & \textbf{90\% CI} \\
\midrule
\textit{E.\ coli}     & 0{,}277 & 0{,}277 & [0{,}159, 0{,}479] \\
\textit{S.\ cerevisiae} & 0{,}240 & 0{,}240 & [0{,}139, 0{,}415] \\
Mamalia              & 0{,}164 & 0{,}164 & [0{,}094, 0{,}285] \\
\bottomrule
\end{tabular}
\end{table}

Seluruh interval kepercayaan berada di bawah 1{,}0, mengkonfirmasi bahwa
pembalikan verdict bersifat robust terhadap ketidakpastian parameter.
Flip $u_G \approx \rho$ menunjukkan bahwa \textbf{begitu fraksi idle
glikolitik melebihi rasio kapasitas, verdict membalik} --- persis kondisi
yang dilaporkan Shen (ekspresi berlebih konstitutif).

\begin{figure}[t]
    \centering
    \includegraphics[width=\columnwidth]{F2_verdict_flip.png}
    \caption{E2: Efisiensi terealisasi vs.\ $u_G$. Garis merah putus-putus = titik-balik; area abu-abu = 90\% CI dari bootstrap. Hijau = respirasi menang, kuning = glikolisis menang.}
    \label{fig:e2}
\end{figure}

\subsection{E2d: Dekomposisi Atribusi Enzim}

Gambar~\ref{fig:decomp} menunjukkan bahwa reatribusi $\pm 30\%$ massa
enzim antar jalur (``enzim mana yang dihitung'') memang menggeser $\rho$
dan titik-balik, namun flip selalu mengikuti $\rho$ secara tepat.
Analisis Sobol mengkonfirmasi bahwa $u_G$ mendominasi ($S_T \approx 0{,}76$--$0{,}78$)
sementara $V_G$ dan $V_R$ masing-masing hanya berkontribusi $S_T \approx 0{,}12$.

\begin{figure}[t]
    \centering
    \includegraphics[width=\columnwidth]{F6_decomposition.png}
    \caption{E2d: Dekomposisi atribusi. Garis biru = flip $u_G$ terhadap faktor reatribusi enzim; garis merah putus = $\rho$ (selalu mengikuti flip). Efek reatribusi nyata namun sekunder.}
    \label{fig:decomp}
\end{figure}

\subsection{E3: Pengukuran Penuntas (Sobol/Morris)}

Gambar~\ref{fig:sobol} menampilkan indeks Sobol total-order ($S_T$)
untuk keenam parameter pada keluaran margin $m$. Untuk ketiga organisme,
$u_G$ mendominasi secara konsisten:

\begin{table}[h]
\centering
\caption{Indeks Sobol total-order ($S_T$) tiga parameter teratas (keluaran: margin $m$).}
\label{tab:sobol}
\begin{tabular}{llcc}
\toprule
\textbf{Organisme} & \textbf{Parameter} & $\boldsymbol{S_T}$ & $\boldsymbol{\pm}$ \\
\midrule
\textit{E.\ coli}    & $u_G$        & 0{,}756 & 0{,}059 \\
                      & $\gamma_G$   & 0{,}122 & 0{,}014 \\
                      & $V_G$        & 0{,}122 & 0{,}014 \\
\midrule
\textit{S.\ cerevisiae} & $u_G$     & 0{,}765 & 0{,}059 \\
                      & $\gamma_G$   & 0{,}124 & 0{,}014 \\
                      & $V_G$        & 0{,}124 & 0{,}014 \\
\midrule
Mamalia               & $u_G$        & 0{,}781 & 0{,}059 \\
                      & $\gamma_G$   & 0{,}126 & 0{,}014 \\
                      & $V_G$        & 0{,}126 & 0{,}015 \\
\bottomrule
\end{tabular}
\end{table}

Morris $\mu^*$ mengkonfirmasi peringkat yang sama (Gambar~\ref{fig:morris}).
\textbf{Temuan E3:} Pengukuran yang paling mengurangi ketidakpastian
verdict adalah $u_G$ --- fraksi pemanfaatan kapasitas glikolitik.

Namun, karena margin $m$ linier di $u_G$, dominasi ini berpotensi
tautologis. E3b menguji robustness temuan ini.

\begin{figure}[t]
    \centering
    \includegraphics[width=\columnwidth]{F3_sobol_tornado.png}
    \caption{E3: Indeks Sobol total-order ($S_T$) per parameter pada keluaran margin (linier). $u_G$ mendominasi di ketiga organisme ($S_T > 0{,}75$).}
    \label{fig:sobol}
\end{figure}

\begin{figure}[t]
    \centering
    \includegraphics[width=\columnwidth]{F3b_morris_scatter.png}
    \caption{Pemeriksaan silang Morris: $\mu^*$ vs.\ $\sigma$. $u_G$ mendominasi, konsisten dengan Sobol.}
    \label{fig:morris}
\end{figure}

\subsection{E3b: Validasi Nonlinier --- Sobol pada Keluaran LP}

Gambar~\ref{fig:sobol_lp} membandingkan indeks Sobol dari keluaran
margin (linier) versus keluaran LP $f^*_G$ (nonlinier).
Pada keluaran nonlinier, kepentingan terdistribusi lebih merata:

\begin{table}[h]
\centering
\caption{Perbandingan $S_T$ teratas: margin (linier) vs.\ LP $f^*_G$ (nonlinier).}
\label{tab:sobol_lp}
\begin{tabular}{llcc}
\toprule
\textbf{Organisme} & \textbf{Parameter} & $\boldsymbol{S_T}$ \textbf{margin} & $\boldsymbol{S_T}$ \textbf{LP} \\
\midrule
\textit{E.\ coli}    & $u_G$        & 0{,}756 & 0{,}633 \\
                      & $\gamma_R$   & 0{,}017 & 0{,}301 \\
\midrule
\textit{S.\ cerevisiae} & $u_G$     & 0{,}765 & 0{,}606 \\
                      & $\gamma_R$   & 0{,}017 & 0{,}343 \\
\midrule
Mamalia               & $\gamma_R$   & 0{,}017 & \textbf{0{,}429} \\
                      & $u_G$        & 0{,}781 & 0{,}395 \\
\bottomrule
\end{tabular}
\end{table}

\textbf{Temuan E3b:} Pada keluaran nonlinier, $u_G$ tetap mendominasi
untuk \textit{E.\ coli} dan ragi, namun untuk sel mamalia
\textbf{$\gamma_R$ mengambil alih peringkat pertama} ($S_T = 0{,}43$
vs.\ $0{,}40$). Ini menunjukkan bahwa (i)~dominasi $u_G$ pada E3 bukan
sepenuhnya tautologis --- ia memang penting secara nonlinier, namun
(ii)~yield ATP respirasi ($\gamma_R$) juga merupakan pengukuran
penuntas yang kritis, terutama untuk sel mamalia. Temuan ini tidak
terlihat dari analisis margin linier.

\begin{figure}[t]
    \centering
    \includegraphics[width=\columnwidth]{F3c_sobol_lp.png}
    \caption{E3b: Perbandingan $S_T$ antara keluaran margin (linier, bar terang) dan keluaran LP $f^*_G$ (nonlinier, bar gelap). Pada keluaran nonlinier, $\gamma_R$ muncul sebagai parameter penting kedua.}
    \label{fig:sobol_lp}
\end{figure}

\subsection{E4: Diagram Fase Metabolik}

Gambar~\ref{fig:phase} menampilkan peta rezim 2D pada sumbu
$(g_\text{avail}, u_G)$ untuk ketiga organisme. Tiga rezim terlihat
jelas: (i)~respirasi dominan (teal, kiri bawah --- glukosa rendah),
(ii)~glikolisis dominan (kuning, kanan atas --- glukosa tinggi dan
$u_G$ tinggi), dan (iii)~zona transisi sempit. Kontur $f^*_G = 0{,}5$
(garis hitam solid) menandai batas keputusan utama.

\textbf{Temuan E4:} Batas rezim bergeser ke kanan (glukosa lebih tinggi)
seiring $u_G$ menurun --- artinya sel dengan enzim glikolitik idle
memerlukan glukosa lebih banyak sebelum beralih ke fermentasi. Ini
konsisten dengan mekanisme \textit{proteome hedging} Shen: sel
mempertahankan enzim glikolitik idle sebagai cadangan, sehingga
ambang fermentasi naik.

\begin{figure}[t]
    \centering
    \includegraphics[width=\columnwidth]{F7_phase_diagram.png}
    \caption{E4: Diagram fase metabolik. Heatmap menunjukkan fraksi fermentatif pada grid $(g_\text{avail}, u_G)$. Kontur hitam = batas $f^*_G = 0{,}5$; putus-putus = $0{,}05$ dan $0{,}95$.}
    \label{fig:phase}
\end{figure}
